
\usepackage{tabularx}
\usepackage{plantuml}
\usepackage{lipsum}  % For dummy text
\usepackage{float}
% ----- Encoding & Font -----
\usepackage[utf8]{inputenc}
\usepackage[T1]{fontenc}

% ----- Cỡ chữ 13pt -----
%
% Kế hoạch môn học yêu cầu Times New Roman 13pt. Tuỳ chọn của \documentclass
% chuẩn chỉ nhận 10/11/12pt, không có 13pt, nên phải dùng gói `fontsize` để đặt
% cỡ chữ nền tuỳ ý mà vẫn giữ nguyên lớp tài liệu `report`. Cách khác là đổi
% sang lớp KOMA-Script (scrreprt hỗ trợ fontsize=13pt), nhưng đổi lớp sẽ kéo
% theo thay đổi cách trình bày tiêu đề chương và mục lục — thay đổi rộng hơn
% nhiều so với thứ đang cần.
\usepackage[fontsize=13pt]{fontsize}

% --- Colors (bắt buộc trước listings) ---
\usepackage[dvipsnames,table]{xcolor}
% dvipsnames: thêm nhiều tên màu có sẵn; 'table' để màu trong bảng booktabs nếu cần.

% ----- Vietnamese? -----
% Nếu bạn cần tiếng Việt, thêm:
\usepackage[vietnamese, english]{babel}

% ----- Margin -----
\usepackage[a4paper, margin=2.5cm]{geometry}

% ----- Line spacing ~1.5 -----
\usepackage{setspace}
\onehalfspacing

% ----- Graphics -----
\usepackage{graphicx}  % Chèn hình
\graphicspath{{../images/}}


% ----- Caption formatting -----
\usepackage[font=small]{caption}


% ----- Table packages -----
\usepackage{array}
\usepackage{booktabs}
\usepackage{multirow}
\usepackage{longtable}  % Bảng trải dài nhiều trang (danh mục viết tắt)

% ----- Math + Times New Roman -----
%
% newtxtext/newtxmath là bản clone Times chất lượng cao dùng được với pdflatex
% (engine đang dùng). Không dùng fontspec + font Times New Roman thật vì
% fontspec đòi XeLaTeX hoặc LuaLaTeX, tức đổi cả toolchain biên dịch.
%
% Thứ tự nạp là bắt buộc: newtxmath phải đứng SAU amsmath. Và không nạp kèm
% amssymb — newtxmath đã cung cấp bộ ký hiệu AMS, nạp cả hai sẽ báo lỗi trùng
% định nghĩa lệnh. Đã kiểm tra toàn bộ chương: không chỗ nào dùng ký hiệu riêng
% của amssymb, chỉ có \times là lệnh LaTeX cơ bản.
\usepackage{amsmath}
\usepackage{newtxtext,newtxmath}

% ----- Hyperlinks -----
\usepackage[hidelinks]{hyperref}

% ----- Code -----
\usepackage{listings}
\lstset{
  basicstyle=\ttfamily\small,
  frame=none,
  breaklines=true,
  breakatwhitespace=true,
  columns=fullflexible,
  keepspaces=true,
  captionpos=b,
  xleftmargin=0.5cm,
  xrightmargin=0.5cm
}

\lstdefinelanguage{JavaScript}{
  keywords={break, case, catch, class, const, continue, debugger, default,
    delete, do, else, export, extends, finally, for, function, if, import,
    in, instanceof, let, new, return, super, switch, this, throw, try,
    typeof, var, void, while, with, yield, async, await},
  keywordstyle=\color{blue}\bfseries,
  ndkeywords={true,false,null,undefined},
  ndkeywordstyle=\color{magenta},
  identifierstyle=\color{black},
  sensitive=true,
  comment=[l]{//},
  morecomment=[s]{/*}{*/},
  commentstyle=\color{gray}\ttfamily,
  stringstyle=\color{red}\ttfamily,
  morestring=[b]',
  morestring=[b]"
}


% ----- Bibliography -----
\usepackage{cite}
% ----- Chống tràn lề -----
%
% Tài liệu có nhiều đường dẫn dài đặt trong \texttt (ví dụ
% docs/results/node-load-test.json). Chữ máy chữ mặc định không được ngắt âm
% tiết, nên khi một đường dẫn không vừa dòng, TeX để nó tràn hẳn ra ngoài lề
% thay vì xuống dòng — phần chữ thừa bị cắt mất khi in.
%
% Hai biện pháp bổ sung cho nhau:
%
% 1. hyphenat với tuỳ chọn htt cho phép ngắt âm tiết trong chữ máy chữ, nhờ đó
%    đường dẫn dài có thể xuống dòng.
% 2. emergencystretch cho TeX thêm biên co giãn khoảng trắng ở lượt sắp chữ
%    cuối cùng, xử lý những dòng còn lại mà việc ngắt âm tiết không cứu được.
%    Giá phải trả là vài dòng có khoảng cách chữ rộng hơn bình thường — chấp
%    nhận được, vì thay thế cho nó là chữ tràn ra ngoài trang giấy.
\setlength{\emergencystretch}{3em}

% ----- Ngắt dòng cho URL dài -----
%
% Thư mục tham khảo có vài URL rất dài (bài blog kỹ thuật của Netflix), và
% \url mặc định chỉ ngắt được ở một số ký tự nhất định. Khi không tìm được chỗ
% ngắt, cả URL tràn ra ngoài lề phải và phần đuôi bị cắt mất khi in. Gói xurl
% cho phép ngắt ở bất kỳ vị trí nào trong URL, đổi một dòng xấu lấy một dòng
% đọc được.
%
% Phải nạp SAU hyperref vì nó chỉnh lại cách hyperref sắp chữ URL.
\usepackage{xurl}
